Comparando tres clasificadores supervisados ---KNN, árbol de decisión y SVM---
sobre dos conjuntos de datos biomédicos de complejidad muy distinta, se
contrastaron sus principios teóricos con el desempeño observado en la
práctica.

El conjunto \textbf{Cancer Wisconsin} resultó ser un problema casi
linealmente separable: incluso en la proyección 2D de PCA las clases benigna
y maligna quedan claramente delimitadas, y los tres modelos alcanzan
exactitudes superiores a $0.94$ sobre el conjunto de prueba. El SVM lineal
fue el más preciso ($0.9766$), capturando la frontera con un margen amplio
y un solo falso negativo de 64 casos malignos.

\textbf{Heart Disease}, en cambio, planteó un escenario mucho más realista
y exigente: $66\%$ de faltantes en \texttt{ca}, $53\%$ en \texttt{thal}, 13
predictoras heterogéneas (continuas y categóricas codificadas) y un
solapamiento marcado entre sanos y enfermos en el espacio reducido. Aquí
ningún modelo superó el $0.85$ de exactitud y el ranking se invirtió: KNN
quedó al frente ($0.8442$), seguido del SVM-RBF ($0.8225$) y muy por
debajo el árbol de decisión ($0.7536$).

PCA cumplió papeles distintos en cada conjunto. En Cancer el método MLE
seleccionó 29 de 30 componentes, indicando que las características son
prácticamente todas independientes y poco redundantes; el umbral del $95\%$
se logra con 10 componentes y la separación entre clases es visible ya en
PC1-PC2. En Heart Disease, MLE escogió 11 de 13 componentes y el $95\%$ se
alcanza con 12: hay incluso menos redundancia, y la dificultad del problema
no se debe a la dimensionalidad sino a la geometría intrínseca de las
clases.

La aplicación de SMOTE, aunque pedida por el enunciado, era prescindible
en ambos casos por el desbalance moderado de partida ($0.59$ y $1.24$).
Su efecto colateral fue introducir ruido sintético del que el KNN se
benefició en Heart Disease (al multiplicar la densidad local de la clase
minoritaria) y posiblemente perjudicó al KNN en Cancer (memorización
excesiva, evidenciada por la brecha CV-test de $0.033$).

El análisis de hiperparámetros mostró que \textbf{$C$ es el parámetro más
    sensible del SVM} en ambos conjuntos, confirmando su rol estructural en
la geometría del margen. Para KNN, \texttt{weights} dominó en el caso
fácil (Cancer) y \texttt{n\_neighbors} en el difícil (Heart Disease):
cuando las clases están bien separadas la métrica de ponderación importa
poco; cuando se solapan, el tamaño del vecindario decide si el modelo
memoriza ruido o suaviza demasiado. El árbol cambió de \texttt{criterion} a
\texttt{max\_depth} entre conjuntos, lo que es coherente con su naturaleza
recursiva: en problemas fáciles cualquier criterio converge a un árbol
similar, pero en problemas difíciles la profundidad es lo único que separa
el sobreajuste del subajuste.

Finalmente, la diferencia de desempeño entre Cancer y Heart Disease, mucho
mayor que cualquier diferencia entre modelos dentro de un mismo conjunto,
deja una conclusión transversal: \textbf{la dificultad intrínseca del
    problema gobierna el techo de exactitud, mientras que la elección del
    clasificador refina sólo los últimos puntos porcentuales}. En problemas
clínicos como este, invertir en calidad de datos, ingeniería de
características y manejo cuidadoso de los faltantes rinde tanto o más que
elegir el algoritmo de última generación.

\QED
